\section{Scope and contributions}

GLEE evaluates agents in bargaining, bilateral trade, and persuasion with
optional natural-language interaction, heterogeneous horizons, and incomplete
information \citep{shapira2024glee}. We study a deliberately inspectable
entrant: its economic actions are produced by fixed Python rules rather than
model inference. This makes failures attributable to explicit state summaries,
branch logic, validation, or the run controller.

This paper makes three bounded contributions. First, it gives a self-contained
specification of the six deployed role policies and their memory boundaries
(Appendix~\ref{app:policy}). Second, it audits failures that legal-action rate or
aggregate rating alone would hide: cycling, uncontrolled exposure, role
reversals, and a monitoring-order bug. Third, it releases the exact executable,
ordered evidence, analysis code, and a 248-game sequential trace. We do
\emph{not} claim strategic superiority, causal mechanism effects, dynamic
receiver obedience, or generalization beyond this run.

The methodological novelty is the alignment of three normally separate audit
layers---strategic decisions, contract enforcement, and sequential exposure---
on one immutable event trace. This decomposition reveals when legal execution,
apparently safe monitoring, and rating success disagree; it is an audit design,
not a new equilibrium or learning algorithm.
